% ============================================================
% Relatório: análise e redução de tráfego do FACEIT (CS2) com Wireshark
% Compilar com: pdflatex relatorio_faceit.tex  (duas vezes)
% Opcional: coloca ao lado deste ficheiro as capturas de ecrã
%   fig_sni.png, fig_dns.png, fig_udp.png  (aparecem automaticamente)
% EDITAR: autores, unidade curricular, grupo, data.
% ============================================================
\documentclass[a4paper,11pt]{article}

\usepackage[utf8]{inputenc}
\usepackage[T1]{fontenc}
\usepackage{lmodern}
\usepackage[portuguese]{babel}
\usepackage[margin=2.5cm]{geometry}
\usepackage{amsmath}
\usepackage{booktabs,array,tabularx,graphicx,caption,microtype}
\usepackage{xcolor}
\usepackage{listings}
\usepackage{pgfplots}
\pgfplotsset{compat=1.16}
\usepackage[hidelinks]{hyperref}

\captionsetup{font=small,labelfont=bf}
\setlength{\parskip}{0.4em}
\emergencystretch=3em
\urlstyle{tt}
\renewcommand{\arraystretch}{1.15}

\lstset{
  basicstyle=\ttfamily\footnotesize,
  breaklines=true,
  frame=single,
  columns=fullflexible,
  keepspaces=true,
  showstringspaces=false,
  backgroundcolor=\color{black!3}
}

% Figura de captura de ecrã: usa o ficheiro se existir, senão mostra uma caixa
\newcommand{\figecra}[1]{%
  \IfFileExists{#1}%
    {\includegraphics[width=\linewidth]{#1}}%
    {\fbox{\parbox{0.92\linewidth}{\centering\vspace{1.4cm}%
      \texttt{\detokenize{#1}}\ (captura de ecrã por inserir)\vspace{1.4cm}}}}}

\title{Análise e redução de volume de tráfego do FACEIT (CS2)\\[0.3em]
\large Captura global com Wireshark, classificação, filtragem e agregação em fluxos}
\author{[Nome do aluno 1] ([n.º]) \and [Nome do aluno 2] ([n.º])\\[0.3em]
[Unidade curricular] --- [Curso / Instituição]}
\date{\today}

\begin{document}
\maketitle

\begin{abstract}
\noindent
Capturou-se, com o Wireshark, todo o tráfego de uma máquina durante uma partida de CS2 na plataforma FACEIT
(326\,122 pacotes, 206,9~MB, cerca de 30~minutos). O serviço foi classificado em dois planos: o \emph{plano de jogo}
(um único fluxo UDP para \texttt{141.95.32.214:27031}, 62\,\% dos pacotes e 79\,\% dos bytes) e o
\emph{plano de controlo} (DNS, TLS e QUIC para \texttt{193.41.200.101/102}, cerca de 1\,\% do volume).
Partindo da mesma captura global, quantificou-se a redução de volume obtida por filtragem
(21,3\,\% dos bytes), truncagem de pacotes com \emph{snaplen} 96 (87,5\,\%), \emph{sampling} 1/10
combinado com \emph{snaplen} (98,9\,\%) e agregação em fluxos com o Argus (99,5\,\% a 99,97\,\%).
\end{abstract}

%------------------------------------------------------------
\section{Introdução e objetivos}
Este trabalho consolida as etapas estudadas sobre captura e análise de tráfego. Os objetivos são:
\begin{itemize}
  \item escolher um serviço (aqui, o FACEIT, com partidas de CS2), gerar tráfego e capturar o tráfego global com o Wireshark;
  \item estudar os critérios de classificação e filtragem mais eficazes (filtros de captura e de visualização);
  \item reduzir o volume de dados capturado/armazenado e quantificar a percentagem de otimização face à captura global,
        por filtragem, truncagem de pacotes, \emph{sampling} e agregação;
  \item explorar uma ferramenta de agregação de pacotes em fluxos a partir de ficheiros \texttt{pcap}.
\end{itemize}

%------------------------------------------------------------
\section{Metodologia}
\paragraph{Captura.} O Wireshark correu na própria máquina de jogo (\texttt{192.168.1.82}), em ligação por cabo,
\textbf{sem filtro de captura e com pacotes inteiros} (captura global), durante 1\,818,6~s (cerca de 30~minutos),
com o FACEIT Anti-Cheat ativo. Antes da captura final fez-se um ensaio num servidor normal de CS2 (fora do FACEIT)
para validar o procedimento.

\paragraph{Estratégia de análise.} Fez-se uma única captura global, convertida para \texttt{pcap} simples
(\texttt{editcap -F pcap}) para servir de \emph{baseline}. Todos os cenários reduzidos foram \textbf{derivados offline
desse mesmo ficheiro}, para que a comparação use exatamente o mesmo tráfego. Usaram-se o Wireshark/\texttt{tshark},
\texttt{editcap}, \texttt{capinfos}, \texttt{tcpdump} e o Argus (agregação em fluxos).

\paragraph{Métrica.} A percentagem de otimização é
\[
  \text{redução} = \Bigl(1 - \frac{\text{tamanho reduzido}}{\text{tamanho global}}\Bigr)\times 100,
\]
medida em bytes (tamanho do ficheiro) e em número de pacotes. O tamanho do ficheiro inclui 16~B de cabeçalho por pacote
no formato \texttt{pcap}; o campo \emph{Data size} do \texttt{capinfos} não reflete a truncagem, por isso não foi usado.

%------------------------------------------------------------
\section{Caracterização do tráfego capturado}
A captura global tem 326\,122 pacotes, 201,6~MB de dados (212,6~MB em \texttt{pcapng}, 206,9~MB em \texttt{pcap}),
com média de 618~B por pacote e 179 pacotes/s. A Tabela~\ref{tab:comp} mostra a composição por classe
(as classes sobrepõem-se parcialmente, pelo que não se somam).

\begin{table}[htbp]
\centering
\caption{Composição do tráfego global.}\label{tab:comp}
\begin{tabularx}{\linewidth}{@{}Xrrrr@{}}
\toprule
Classe & Pacotes & Bytes & \% pacotes & \% bytes\\
\midrule
Servidor de jogo (\texttt{141.95.32.214})      & 202\,086 & 159\,679\,469 & 62,0 & 79,2\\
UDP \texttt{:19336} (\texttt{104.29.146.69})   &  85\,785 &  17\,938\,347 & 26,3 &  8,9\\
TCP (todo)                                      &  19\,394 &  12\,376\,258 &  5,9 &  6,1\\
QUIC (todo)                                     &   7\,226 &   9\,622\,529 &  2,2 &  4,8\\
mDNS / SSDP / DNS                               &   4\,576 &   1\,307\,245 &  1,4 &  0,6\\
Controlo FACEIT (\texttt{193.41.200.101/102})   &   3\,995 &   3\,338\,873 &  1,2 &  1,7\\
\bottomrule
\end{tabularx}
\end{table}

\paragraph{Plano de jogo.} A conversa UDP \texttt{192.168.1.82:50088}~$\leftrightarrow$~\texttt{141.95.32.214:27031}
domina a captura (Figura~\ref{fig:udp}): começa aos 293,8~s, dura 1\,520,5~s e tem 202\,006 pacotes e 159,7~MB.
O cliente envia 96\,786 pacotes (70,7~MB, média de 730~B) e o servidor 105\,220 (89,0~MB, média de 846~B).
O débito é praticamente constante: cerca de 133 pacotes/s e 0,84~Mbit/s (aproximadamente 40\,000 pacotes por
cada 300~s). Os tamanhos concentram-se entre 640 e 1\,279~B (69,3\,\% dos pacotes), com média de 790~B.

\paragraph{Plano de controlo.} O DNS e o SNI (TLS e QUIC) identificam o serviço: \nolinkurl{www.faceit.com},
\nolinkurl{edge.faceit.com}, \nolinkurl{chat.faceit.com}, \nolinkurl{faceit.com}, \nolinkurl{cdn-frontend.faceit-cdn.net},
\nolinkurl{distribution.faceit-cdn.net} e \nolinkurl{ac-client-ws.faceit.com}, todos resolvidos para
\texttt{193.41.200.101} e \texttt{193.41.200.102} (Figuras~\ref{fig:sni} e~\ref{fig:dns}), mais
\texttt{analytics.faceitanalytics.com} (Cloudflare). O primeiro contacto ocorre aos 2,2~s (QUIC para
\texttt{www.faceit.com}), há nova rajada de pedidos aos 225--245~s, antes de o fluxo de jogo começar, e um contacto
a \texttt{ac-client-ws.faceit.com} aos 1\,804,6~s. Este plano é pequeno: 3\,995 pacotes e 3,3~MB.

\paragraph{Ruído.} Os fluxos UDP para \texttt{104.29.146.69:19336} (um novo fluxo por minuto, 85\,785 pacotes) são
independentes do jogo: apareciam também na captura de ensaio. A sua origem não foi identificada. O restante é tráfego
da rede local (mDNS, SSDP, \emph{broadcasts}) e outras ligações HTTPS/QUIC.

\begin{figure}[htbp]
\centering
\figecra{fig_udp.png}
\caption{\emph{Conversations}~$\rightarrow$~UDP, ordenado por pacotes: o fluxo para \texttt{141.95.32.214:27031} domina.}\label{fig:udp}
\end{figure}

\begin{figure}[htbp]
\centering
\figecra{fig_sni.png}
\caption{Filtro \texttt{tls.handshake.extensions\_\allowbreak server\_name contains "faceit"} (TLS e QUIC).}\label{fig:sni}
\end{figure}

\begin{figure}[htbp]
\centering
\figecra{fig_dns.png}
\caption{Filtro \texttt{dns.qry.\allowbreak name contains "faceit"}.}\label{fig:dns}
\end{figure}

%------------------------------------------------------------
\section{Critérios de classificação e filtros}
A Tabela~\ref{tab:filtros} resume os critérios estudados. O filtro de captura (BPF) só vê cabeçalhos
(IP, protocolo, portas), pelo que critérios baseados em nomes (DNS/SNI) só são possíveis como filtros de visualização.
A eficácia mede-se por cobertura (pacotes do jogo apanhados / pacotes do jogo) e precisão (pacotes do jogo
apanhados / pacotes apanhados), tomando como referência a conversa do jogo (202\,006 pacotes).

\begin{table}[htbp]
\centering
\caption{Critérios de classificação e eficácia (BPF = filtro de captura; Vis. = filtro de visualização).}\label{tab:filtros}
\small
\begin{tabularx}{\linewidth}{@{}p{2.7cm}XXrr@{}}
\toprule
Critério & BPF & Vis. & Pacotes & Precisão\\
\midrule
Servidor (IP)     & \texttt{udp and host 141.95.32.214}   & \texttt{ip.addr==141.95.32.214 \&\& udp}    & 202\,006 & 100\,\%\\
Qualquer protocolo com o IP & \texttt{host 141.95.32.214}  & \texttt{ip.addr==141.95.32.214}             & 202\,060 & 99,97\,\%\\
Gama de portas    & \texttt{udp portrange 27000-27050}    & \texttt{udp.port>=27000 \&\& udp.port<=27050} & 202\,040 & 99,98\,\%\\
Só UDP            & \texttt{udp}                          & \texttt{udp}                                & 300\,571 & 67,2\,\%\\
Controlo (IP)     & \texttt{host 193.41.200.101 or host 193.41.200.102} & \texttt{ip.addr==193.41.200.101 \textbar\textbar{} ip.addr==193.41.200.102} & 3\,995 & ---\\
Controlo (nome)   & (não aplicável) & \texttt{dns.qry.\allowbreak name contains "faceit"}; \texttt{tls.handshake.extensions\_\allowbreak server\_name contains "faceit"} & --- & ---\\
Jogo + controlo   & \texttt{host} servidor \texttt{or host} controlo (2 IPs) & --- & 206\,055 & ---\\
\bottomrule
\end{tabularx}
\end{table}

Todos os filtros por IP ou por portas têm cobertura de 100\,\% do fluxo de jogo. A gama de portas apanha apenas 34 pacotes
a mais do que o filtro por IP (precisão de 99,98\,\%) e tem a vantagem de não depender do IP do servidor, que muda de partida
para partida. O filtro apenas por protocolo (\texttt{udp}) é demasiado largo: 33\,\% do que apanha não é jogo.
Como critérios de deteção do fluxo de jogo funcionam bem a combinação de UDP, a gama de portas \texttt{27xxx} e a
dominância por volume (\emph{Conversations} ordenado por bytes), confirmada pelo DNS/SNI do plano de controlo.

%------------------------------------------------------------
\section{Redução de volume}
A Tabela~\ref{tab:red} e a Figura~\ref{fig:red} apresentam os cenários. Todas as percentagens são cumulativas face ao
\emph{baseline} (A) de 326\,122 pacotes e 206\,864\,624~B.

\begin{table}[htbp]
\centering
\caption{Redução de volume face à captura global (percentagens em tamanho de ficheiro e em pacotes).}\label{tab:red}
\small
\begin{tabularx}{\linewidth}{@{}lXrrrr@{}}
\toprule
 & Cenário & Pacotes & Bytes & Red. bytes & Red. pacotes\\
\midrule
A  & Captura global (\emph{baseline})              & 326\,122 & 206\,864\,624 & 0\,\%     & 0\,\%\\
B1 & Filtro \texttt{udp and host} servidor         & 202\,006 & 162\,905\,309 & 21,25\,\% & 38,06\,\%\\
B3 & Filtro gama de portas 27000--27050            & 202\,040 & 162\,951\,481 & 21,23\,\% & 38,05\,\%\\
B4 & Filtro só \texttt{udp}                        & 300\,571 & 193\,613\,964 &  6,41\,\% &  7,83\,\%\\
B5 & Filtro jogo + controlo FACEIT                 & 206\,055 & 166\,312\,206 & 19,60\,\% & 36,82\,\%\\
C1 & B1 + \emph{snaplen} 128                       & 202\,006 &  32\,311\,732 & 84,38\,\% & 38,06\,\%\\
C2 & B1 + \emph{snaplen} 96                        & 202\,006 &  25\,856\,584 & 87,50\,\% & 38,06\,\%\\
C3 & B1 + \emph{snaplen} 64                        & 202\,006 &  19\,392\,644 & 90,63\,\% & 38,06\,\%\\
C4 & B5 + \emph{snaplen} 96                        & 206\,055 &  26\,333\,772 & 87,27\,\% & 36,82\,\%\\
D1 & B1 + \emph{sampling} 1/10                     &  20\,200 &  16\,376\,628 & 92,08\,\% & 93,81\,\%\\
D2 & B1 + \emph{snaplen} 96 + \emph{sampling} 1/10 &  20\,200 &   2\,262\,424 & 98,91\,\% & 93,81\,\%\\
E1 & Fluxos Argus, global (registos de 5~s)        & --- &   1\,004\,872 & 99,51\,\% & ---\\
E2 & Fluxos Argus, global (1 registo/fluxo, CSV)   & --- &      64\,839 & 99,97\,\% & ---\\
E3 & Fluxos Argus, só servidor (CSV, 1 fluxo)      & --- &         229 & $\approx$100\,\% & ---\\
\bottomrule
\end{tabularx}
\end{table}

\begin{figure}[htbp]
\centering
\begin{tikzpicture}
\begin{axis}[
  xbar, width=0.74\linewidth, height=9cm,
  xmin=0, xmax=115, xlabel={Redução face à captura global (\% dos bytes)},
  symbolic y coords={B1 filtro IP, B3 filtro portas, B5 jogo+controlo, C1 snaplen 128, C2 snaplen 96, C3 snaplen 64, D1 sampling 1/10, D2 snaplen 96 + sampling, E1 fluxos 5 s, E2 fluxos 1/fluxo},
  ytick=data, y dir=reverse, enlarge y limits=0.07, bar width=7pt,
  yticklabel style={font=\footnotesize}, xmajorgrids,
  nodes near coords={\pgfmathprintnumber[use comma,fixed,fixed zerofill,precision=2]{\pgfplotspointmeta}},
  every node near coord/.append style={font=\scriptsize}
]
\addplot[fill=blue!55,draw=blue!70!black] coordinates {
  (21.25,B1 filtro IP) (21.23,B3 filtro portas) (19.60,B5 jogo+controlo)
  (84.38,C1 snaplen 128) (87.50,C2 snaplen 96) (90.63,C3 snaplen 64)
  (92.08,D1 sampling 1/10) (98.91,D2 snaplen 96 + sampling)
  (99.51,E1 fluxos 5 s) (99.97,E2 fluxos 1/fluxo)};
\end{axis}
\end{tikzpicture}
\caption{Redução de bytes por cenário.}\label{fig:red}
\end{figure}

\paragraph{Filtrar.} Filtrar o fluxo de jogo reduz 21,25\,\% dos bytes mas 38,06\,\% dos pacotes: o jogo já é 79\,\%
dos bytes capturados, enquanto o ruído é composto por muitos pacotes pequenos (os 85\,785 pacotes de \texttt{:19336}
pesam 26\,\% dos pacotes mas 9\,\% dos bytes). Incluir o plano de controlo no filtro custa pouco
(19,60\,\% contra 21,25\,\%) e descreve o serviço de forma mais completa.

\paragraph{Não capturar o pacote inteiro.} É a técnica com maior impacto individual. Os cabeçalhos Ethernet, IPv4 e UDP
ocupam 42~B, por isso com \emph{snaplen} 96 guardam-se os cabeçalhos e mais 54~B de payload. Com pacotes de
790~B em média, a redução total passa de 21,25\,\% para 87,50\,\% (84,38\,\% com 128~B e 90,63\,\% com 64~B).
O custo é a perda do restante payload.

\paragraph{\emph{Sampling}.} Guardar 1 pacote em 10 (por ordem de pacote) dá 92,08\,\% sem truncagem e 98,91\,\% com
\emph{snaplen} 96 (2,26~MB, dos quais cerca de 0,32~MB são cabeçalhos de pacote do formato \texttt{pcap}). As estatísticas
de volume mantêm-se: extrapolando $\times 10$, o erro é $-0{,}003\,\%$ nos pacotes e $+0{,}54\,\%$ nos bytes.
Perdem-se as sequências de pacotes consecutivos.

%------------------------------------------------------------
\section{Agregação de pacotes em fluxos}
Usou-se o Argus (\texttt{argus -r f.pcap -w f.argus}), que emite registos de estado de cerca de 5~s por fluxo, e o
\texttt{racluster} para obter um registo por fluxo (5-tuplo), exportado para CSV com o \texttt{ra}. Do \emph{baseline}
de 206,9~MB obtêm-se 1,0~MB em registos de 5~s (99,51\,\%), 122,7~kB em binário com um registo por fluxo (99,94\,\%)
e 64,8~kB em CSV (99,97\,\%), com 797 fluxos (mais um registo de gestão do Argus). Aplicado ao filtro do servidor, o fluxo
de jogo fica reduzido a um único registo de 229~B no CSV.

Os totais do fluxo de jogo no Argus (70\,661\,525~B do cliente e 89\,011\,664~B do servidor) somam
159\,673\,189~B, exatamente o volume de dados do pcap filtrado. A agregação conserva o 5-tuplo, o instante de início,
a duração e os pacotes e bytes por sentido, mas perde a temporização e o tamanho de cada pacote e o payload. Os registos de
5~s (E1) conservam ainda a evolução temporal do fluxo.

%------------------------------------------------------------
\section{Conclusões}
\begin{itemize}
  \item O tráfego do FACEIT divide-se em plano de jogo (UDP \texttt{141.95.32.214:27031}, 79\,\% dos bytes) e plano de controlo
        (DNS/TLS/QUIC para \texttt{193.41.200.101/102}, cerca de 1\,\%).
  \item Os critérios mais eficazes para detetar o fluxo de jogo são UDP, gama de portas \texttt{27000--27050} e dominância por
        volume; o DNS/SNI confirma o serviço. O filtro apenas por protocolo é demasiado largo.
  \item Filtrar poupa pouco em bytes (21\,\%) porque o jogo domina a captura. As maiores poupanças vêm do \emph{snaplen}
        (87,5\,\%), do \emph{sampling} (98,9\,\% em conjunto) e da agregação em fluxos (99,5\,\% a 99,97\,\%).
  \item A escolha depende da utilização: cabeçalhos por pacote $\Rightarrow$ filtro e \emph{snaplen} 96; estatísticas de volume
        $\Rightarrow$ \emph{sampling}; contabilização e retenção prolongada $\Rightarrow$ fluxos.
\end{itemize}

\paragraph{Limitações.} Analisou-se uma única captura de uma única partida; o IP e a porta do servidor mudam entre
partidas. Os cenários foram derivados offline da captura global (não capturados ao vivo). O \emph{sampling} é determinístico
(1 em 10 por ordem de pacote). As percentagens em ficheiro incluem o cabeçalho de 16~B por pacote do formato \texttt{pcap}.
Só se usou o Argus para a agregação em fluxos.

%------------------------------------------------------------
\appendix
\section{Comandos utilizados}
\begin{lstlisting}
# Baseline
capinfos faceit.pcapng
editcap -F pcap faceit.pcapng A_global.pcap

# Caracterizacao
tshark -r A_global.pcap -q -z io,phs
tshark -r A_global.pcap -q -z conv,udp
tshark -r A_global.pcap -q -z plen,tree,"ip.addr==141.95.32.214"
tshark -r A_global.pcap -Y 'dns.flags.response==1 && dns.qry.name contains "faceit"' -T fields -e frame.time_relative -e dns.qry.name -e dns.a
tshark -r A_global.pcap -Y 'tls.handshake.extensions_server_name contains "faceit"' -T fields -e frame.time_relative -e ip.dst -e tls.handshake.extensions_server_name

# Filtros (BPF, aplicados offline)
tcpdump -r A_global.pcap -w B1_udp_host.pcap  "udp and host 141.95.32.214"
tcpdump -r A_global.pcap -w B3_portrange.pcap "udp portrange 27000-27050"
tcpdump -r A_global.pcap -w B4_udp.pcap       "udp"
tcpdump -r A_global.pcap -w B5_game_ctrl.pcap "host 141.95.32.214 or host 193.41.200.101 or host 193.41.200.102"

# Snaplen e sampling
editcap -s 96 B1_udp_host.pcap C2_snaplen96.pcap
tshark -r C2_snaplen96.pcap -Y "frame.number % 10 == 0" -F pcap -w D2_snaplen_sampling.pcap

# Fluxos (Argus)
argus -r A_global.pcap -w E1_global.argus
racluster -r E1_global.argus -m proto saddr sport daddr dport -w E2_global_cluster.argus
ra -r E2_global_cluster.argus -n -c , -s stime dur proto saddr sport daddr dport spkts dpkts sbytes dbytes > E2_global_flows.csv

# Medicao
capinfos -T -m -M -c -s *.pcap
\end{lstlisting}

\end{document}
